\documentclass[12pt,a4paper]{article}
\usepackage[utf8]{inputenc}
\usepackage{amsmath,amsfonts,amssymb,amsthm}
\usepackage{geometry}
\geometry{margin=1in}
\usepackage{hyperref}
\usepackage{xcolor}
\usepackage{booktabs}
\usepackage{enumitem}
\usepackage{graphicx}
\usepackage{fancyhdr}

\pagestyle{fancy}
\fancyhf{}
\fancyfoot[L]{Universitas Bengkulu - Teknik Elektro}
\fancyfoot[R]{\thepage}
\def\footrule{{\color{gray}\hrule}}

\theoremstyle{definition}
\newtheorem{definition}{Definisi}[section]
\newtheorem{example}{Contoh}[section]

\title{Modul Tambahan Persamaan Diferensial\\Studi Kasus Industri Modern di Teknik Elektro}
\author{Novalio Daratha\\Program Studi Teknik Elektro\\Universitas Bengkulu}
\date{2026}

\begin{document}

\maketitle

\tableofcontents

\newpage

\section{Pendahuluan}
Modul tambahan ini menyajikan aplikasi persamaan diferensial tingkat lanjut dalam menyelesaikan permasalahan kompleks di industri kelistrikan dan telekomunikasi modern. Pembahasan mencakup stabilitas transien pada sistem tenaga listrik, perambatan sinyal frekuensi tinggi untuk jaringan 5G, serta fenomena \textit{skin effect} pada saluran transmisi tegangan tinggi. Modul ini bertujuan untuk memberikan wawasan aplikatif tentang bagaimana persamaan diferensial digunakan untuk memodelkan dan menyelesaikan masalah nyata di dunia industri teknik elektro.

\section{Studi Kasus 1: Stabilitas Transien Sistem Tenaga Listrik (Persamaan Ayunan)}

\subsection{Latar Belakang Fisika}
Dalam sistem interkoneksi tenaga listrik modern, menjaga sinkronisasi antar generator sangatlah penting, terutama saat terjadi gangguan seperti hubung singkat (korsleting) atau perubahan beban yang mendadak. Fenomena dinamika rotor generator dalam merespons ketidakseimbangan antara daya mekanik dari turbin dan daya elektrik yang disalurkan ke jaringan dimodelkan menggunakan Persamaan Ayunan (\textit{Swing Equation}).

\subsection{Model Persamaan Diferensial}
Persamaan ayunan merupakan Persamaan Diferensial Biasa (PDB) orde 2 non-linear, yang dinyatakan sebagai:
\begin{equation}
    M \frac{d^2 \delta}{dt^2} + D \frac{d \delta}{dt} = P_m - P_e
\end{equation}
dimana:
\begin{itemize}
    \item $M$ adalah konstanta inersia generator.
    \item $\delta$ adalah sudut rotor (dalam radian).
    \item $D$ adalah koefisien redaman (\textit{damping}).
    \item $P_m$ adalah daya mekanik masuk dari turbin.
    \item $P_e$ adalah daya elektrik yang keluar ke jaringan. Biasanya, daya elektrik merupakan fungsi non-linear dari sudut rotor: $P_e = P_{\max} \sin(\delta)$.
\end{itemize}
Sehingga, persamaan lengkapnya menjadi:
\begin{equation}
    M \frac{d^2 \delta}{dt^2} + D \frac{d \delta}{dt} + P_{\max} \sin(\delta) = P_m
\end{equation}

\subsection{Metode Penyelesaian}
Karena adanya suku trigonometri non-linear $\sin(\delta)$, persamaan ini secara umum tidak memiliki solusi analitik tertutup yang sederhana kecuali untuk sudut-sudut kecil (di mana $\sin(\delta) \approx \delta$). 

Untuk studi sistem tenaga di industri nyata, penyelesaian persamaan ini bergantung pada \textbf{metode numerik}. Metode integrasi numerik langkah-demi-langkah seperti \textbf{Metode Euler yang dimodifikasi} atau \textbf{Metode Runge-Kutta orde 4 (RK4)} sering diterapkan menggunakan perangkat lunak komputasi (seperti MATLAB, Octave, atau Julia) untuk mengevaluasi stabilitas sistem sebelum, selama, dan setelah terjadinya gangguan. Evaluasi ini dikenal dengan Kriteria Luas Sama (\textit{Equal Area Criterion}) pada kasus sistem mesin tunggal yang terhubung ke \textit{infinite bus}.

\section{Studi Kasus 2: Perambatan Sinyal Frekuensi Tinggi di Jaringan 5G (Persamaan Telegrafer)}

\subsection{Latar Belakang Fisika}
Pengembangan jaringan 5G menuntut transmisi data dengan laju sangat tinggi (frekuensi dalam orde GHz). Pada frekuensi ini, saluran transmisi tidak dapat lagi direpresentasikan hanya sebagai rangkaian elemen terpusat (\textit{lumped-parameter circuit}), melainkan sebagai rangkaian parameter terdistribusi (\textit{distributed-parameter}). Tegangan dan arus berubah secara kontinu sepanjang saluran akibat fenomena perambatan gelombang elektromagnetik.

\subsection{Model Persamaan Diferensial}
Model matematis saluran transmisi dinyatakan melalui Persamaan Telegrafer (\textit{Telegrapher's Equations}), yang merupakan sistem Persamaan Diferensial Parsial (PDP) tergandeng (\textit{coupled}):
\begin{align}
    -\frac{\partial v(x,t)}{\partial x} &= R i(x,t) + L \frac{\partial i(x,t)}{\partial t} \\
    -\frac{\partial i(x,t)}{\partial x} &= G v(x,t) + C \frac{\partial v(x,t)}{\partial t}
\end{align}
dimana:
\begin{itemize}
    \item $v(x,t)$ adalah tegangan sepanjang posisi $x$ dan waktu $t$.
    \item $i(x,t)$ adalah arus.
    \item $R, L, G, C$ adalah hambatan, induktansi, konduktansi, dan kapasitansi saluran per satuan panjang.
\end{itemize}
Dengan melakukan substitusi dan diferensiasi, sistem ini dapat direduksi menjadi Persamaan Gelombang yang teredam:
\begin{equation}
    \frac{\partial^2 v}{\partial x^2} = L C \frac{\partial^2 v}{\partial t^2} + (R C + L G) \frac{\partial v}{\partial t} + R G v
\end{equation}

\subsection{Metode Penyelesaian}
\textbf{Solusi analitik} untuk kasus gelombang harmonik (kondisi tunak sinusoidal) diperoleh melalui transformasi fasor, mereduksi PDP waktu ke PDB dalam domain frekuensi. Solusi dari persamaan orde dua ini akan menghasilkan dua gelombang: gelombang berjalan (\textit{incident wave}) dan gelombang pantulan (\textit{reflected wave}). Parameter kunci seperti konstanta propagasi ($\gamma = \alpha + j\beta$) dan impedansi karakteristik ($Z_0$) digunakan untuk mendesain \textit{impedance matching} yang krusial di antena 5G guna menghindari rugi-rugi (\textit{return loss}) yang fatal.

\textbf{Solusi numerik}, seperti Finite-Difference Time-Domain (FDTD), digunakan untuk memodelkan struktur saluran mikrostrip yang kompleks dengan berbagai bahan dielektrik di perancangan sirkuit RF (\textit{Radio Frequency}).

\section{Studi Kasus 3: Efek Kulit (\textit{Skin Effect}) pada Saluran Transmisi Tegangan Tinggi (Persamaan Bessel)}

\subsection{Latar Belakang Fisika}
Pada sistem penyaluran energi listrik arus bolak-balik (AC), distribusi arus di dalam konduktor silindris (seperti kawat aluminium bertulang baja - ACSR) tidak merata. Arus cenderung terkonsentrasi di permukaan konduktor pada frekuensi tinggi. Fenomena ini disebut efek kulit (\textit{skin effect}), yang menyebabkan resistansi efektif konduktor meningkat, memperbesar rugi-rugi konduksi (\textit{losses}) seiring meningkatnya dimensi kawat dan frekuensi eksitasi.

\subsection{Model Persamaan Diferensial}
Analisis efek kulit diturunkan dari Persamaan Maxwell. Jika kita mengasumsikan konduktor berbentuk silinder dengan jari-jari $r$ dan arus searah memanjang (sumbu-$z$), kerapatan arus $J_z(r)$ (sebagai fasor) yang bergantung pada jarak dari pusat $r$, memenuhi persamaan Helmholtz dalam koordinat silindris:
\begin{equation}
    \frac{d^2 J_z}{dr^2} + \frac{1}{r} \frac{d J_z}{dr} - j\omega\mu\sigma J_z = 0
\end{equation}
dimana $\omega$ adalah frekuensi sudut, $\mu$ permeabilitas magnetik, dan $\sigma$ konduktivitas listrik.

Misalkan $k^2 = -j\omega\mu\sigma$. Persamaan tersebut menjadi bentuk standar dari Persamaan Diferensial Bessel Orde Nol:
\begin{equation}
    r^2 \frac{d^2 J_z}{dr^2} + r \frac{d J_z}{dr} + k^2 r^2 J_z = 0
\end{equation}

\subsection{Metode Penyelesaian}
\textbf{Penyelesaian Analitik} dari persamaan Bessel tersebut melibatkan fungsi Bessel jenis pertama orde nol, $J_0(kr)$. Mengingat parameter $k$ adalah bilangan kompleks, solusi kerapatan arus diekspresikan dalam bentuk Fungsi Kelvin (\textit{Kelvin functions}, yaitu fungsi \textit{ber} dan \textit{bei}).
\begin{equation}
    J_z(r) = J_z(0) J_0(\sqrt{-j\omega\mu\sigma} r)
\end{equation}
Dari distribusi kerapatan arus ini, resistansi AC per satuan panjang konduktor dapat dievaluasi secara integratif. Di industri modern, insinyur utilitas listrik menggunakan kurva standar yang diturunkan dari solusi Fungsi Bessel ini untuk menghitung kompensasi termal dan \textit{ampacity} (batas kapasitas hantar arus) konduktor pada kondisi beban puncak dengan presisi tinggi tanpa perlu melakukan komputasi numerik secara terus-menerus.

\section{Kesimpulan}
Ketiga studi kasus modern yang telah dijabarkan mendemonstrasikan secara jelas peranan persamaan diferensial di berbagai sektor teknik elektro; dari PDB non-linear pada dinamika kestabilan generator besar, PDP pada propagasi gelombang elektromagnetik berkecepatan tinggi, hingga Persamaan Bessel yang istimewa dalam mengkarakterisasi efek kulit konduktor industri. Keahlian menyelesaikan formulasi matematis ini, baik secara analitik untuk pemahaman dasar maupun secara numerik untuk kompleksitas sejati, merupakan kompetensi krusial bagi calon insinyur elektro tingkat lanjut.

\end{document}
